\chapter{Ingénierie des indicateurs (feature engineering)}
\label{chap:features}

La qualité des indicateurs extraits conditionne directement la performance des modèles. Ce
chapitre détaille la conception des vecteurs de features, dataset par dataset, et leur
justification physique. C'est l'une des contributions les plus déterminantes du projet :
des features bien pensées permettent à des modèles simples d'atteindre d'excellentes
performances.

\section{Features vibratoires standard}
Pour les signaux vibratoires temporels (VBL, CWRU), chaque canal produit \textbf{15
descripteurs} : neuf temporels et six fréquentiels.

\begin{itemize}[leftmargin=1.5em]
  \item \textbf{Temporels (9)} : RMS, kurtosis, skewness, valeur crête, crête-à-crête,
        facteur de crête, facteur de forme, facteur d'impulsion et écart-type.
  \item \textbf{Fréquentiels (6)} : fréquence dominante, entropie spectrale et quatre
        énergies de bande (0--1, 1--3, 3--5 et 5--6{,}4\,kHz).
\end{itemize}

Le choix de ces descripteurs n'est pas arbitraire : le \textbf{kurtosis} et le \textbf{facteur
de crête} sont les premiers à réagir à un défaut de roulement naissant (chocs périodiques),
tandis que le \textbf{RMS} traduit une dégradation déjà installée. Les facteurs de
\emph{forme} et d'\emph{impulsion} complètent la caractérisation de la morphologie du signal.
Côté fréquentiel, les \textbf{quatre bandes d'énergie} localisent l'excitation : les défauts
de roulement privilégient les hautes fréquences (résonances), tandis qu'un balourd reste basse
fréquence. L'entropie spectrale distingue un spectre concentré (raie nette) d'un spectre
diffus (jeu mécanique, bruit large bande). Avec trois canaux, VBL atteint 45 features ; avec
un seul canal, CWRU en compte 15.

\section{Features par cycles (CMAPSS)}
Pour CMAPSS, le signal n'est pas vibratoire mais constitué de relevés de 21 capteurs au fil
des cycles de fonctionnement. Pour chaque capteur, six statistiques sont calculées sur une
fenêtre de cycles : moyenne, écart-type, minimum, maximum, \textbf{pente de régression
linéaire} (tendance) et kurtosis. La pente est cruciale : elle capture la \emph{dérive}
progressive d'un capteur, signal avant-coureur de la dégradation. On obtient ainsi
$21 \times 6 = 126$ features.

\section{Features spécialisées engrenages}
Pour les défauts d'engrenage (MCC5-THU), des descripteurs dédiés ont été conçus, car les
signatures d'engrenage sont difficiles à résumer par des descripteurs génériques :
\begin{itemize}[leftmargin=1.5em]
  \item l'\textbf{énergie autour de la fréquence d'engrènement} (GMF $=$ nombre de dents
        $\times$ fréquence de rotation) et de ses premières harmoniques ;
  \item les \textbf{bandes latérales} (sidebands) à GMF $\pm\,k\!\times$ la fréquence de
        rotation, dont la présence signe une \emph{modulation} caractéristique d'un défaut
        d'engrenage ;
  \item le \textbf{ratio sidebands/GMF}, indicateur synthétique de modulation ;
  \item la \textbf{démodulation d'enveloppe} (Hilbert) sur la bande de résonance ;
  \item l'\textbf{analyse cepstrale}, qui détecte les familles harmoniques périodiques liées
        à la rotation.
\end{itemize}

\section{Features multi-capteurs (Mechanical Faults)}
Le cas le plus élaboré est Mechanical Faults, dont les \textbf{132 features} (33 par canal
$\times$ 4 canaux) ont été conçues pour discriminer balourd, désalignement et jeu mécanique.
Le tableau~\ref{tab:mffeat} détaille la composition par canal.

\begin{table}[h]
\centering\small
\caption{Composition du vecteur de features Mechanical Faults (par canal).}
\label{tab:mffeat}
\begin{tabular}{>{\bfseries}l c p{7.2cm}}
\toprule
Groupe & Nombre & Rôle physique \\
\midrule
Temporels & 6 & RMS, kurtosis, skewness, crête-à-crête, facteur de crête, écart-type \\
Harmoniques rotation & 10 & énergie à $1\times$,$2\times$,$3\times$,$4\times$,$5\times$ et sous-harmoniques ($0{,}5\times$, $1/3\times$, etc.) \\
Ratio énergie & 1 & énergie totale / fondamentale \\
Bandes larges & 6 & répartition d'énergie (seuils 0--50--200--500--1k--3k--6{,}4k\,Hz) \\
Enveloppe Hilbert & 5 & RMS, kurtosis, écart-type, énergie enveloppe à $1\times$ et $2\times$ \\
Ratios harmoniques & 5 & $1\times/2\times$ (balourd), $2\times/3\times$ (désalignement), dominance fondamentale, etc. \\
\midrule
\textbf{Total / canal} & \textbf{33} & $\times 4$ canaux $= 132$ features \\
\bottomrule
\end{tabular}
\end{table}

Les \textbf{ratios harmoniques} encodent directement la \emph{physique} des défauts : un
balourd se caractérise par $E(1\times) \gg E(2\times)$ (ratio élevé) ; un désalignement par
$E(2\times) \approx E(1\times)$ (ratio proche de 1) ; un jeu mécanique par une énergie diffuse
sur de nombreuses harmoniques. En fournissant directement ces ratios au modèle, on lui
transmet une connaissance experte qui explique pourquoi le MLP atteint 89{,}6\,\% là où une
approche naïve échouerait. C'est l'illustration concrète du principe selon lequel
\emph{de bonnes features valent mieux qu'un modèle complexe}.

\section{Cache et reproductibilité}
L'extraction des features étant coûteuse, un mécanisme de \textbf{cache} (clé MD5 calculée sur
le chemin du fichier et la configuration YAML) évite tout recalcul inutile et garantit la
reproductibilité des expériences. Les tenseurs de features ($X$), labels ($y$) et métadonnées
sont sérialisés (\texttt{.npy}, \texttt{.pkl}) pour alimenter directement l'entraînement des
modèles et la plateforme en production.
